\documentclass[11pt,a4paper]{article}
\usepackage[utf8]{inputenc}
\usepackage[top=0.9in, bottom=0.8in, left=1in, right=1in]{geometry}
\usepackage{hyperref}

\setlength{\parskip}{0.8em}
\setlength{\parindent}{0pt}

\begin{document}

\begin{center}
    {\LARGE \textbf{Elias Mei}} \\
    \vspace{2mm}
    Linköping, Sweden $\cdot$ liume102@student.liu.se $\cdot$ +46-0704898962
    \vspace{4mm}
\end{center}

\begin{flushright}
May 15, 2026
\end{flushright}


Dear Dr. Anne-Marie Fors Connolly,

I am writing to apply for the PhD position in Data Driven Life Sciences at Umeå University. I am completing my master's degree in Statistics and Machine Learning at Linköping University, and I will graduate in June this year.

\textbf{Motivation for the position} I believe the most interesting work in AI right now happens at the intersection with other fields, and healthcare is where that intersection feels most urgent. This position is about building machine learning methods that turn messy clinical microbiology data into something structured and usable for epidemiologists and clinicians. That combination of algorithmic thinking and real-world impact is why I moved from software engineering into computational biology.

I had already come across two lines of work from this group before seeing the position. One is the Infomap review by Rosvall and colleagues \cite{smiljanic2025mapequation}, which develops community detection from how information flows through a network. It is essentially graph mining, and the approach resonated with me. In my thesis, I used graph attention networks on OmniPath to predict drug response, which shares a similar philosophy of letting the data decide which connections matter. The other is the group's nationwide cohort study using Swedish registry data \cite{jakobsson2025hepatitis}. What struck me there was the scale, millions of linked records across multiple registries, with all the messiness that implies. That is not a toy dataset, and making sense of it requires serious data engineering. After five years of building large-scale data pipelines in industry, I know exactly what that kind of work demands.

\textbf{Suitability for research topic} I have spent over five years as a software engineer building data pipelines at scale. At SAP, I wrote systems that parsed enterprise text logs and turned them into structured, queryable tables. At NavInfo and Ericsson, I designed distributed microservices that processed streaming data in real time. I am comfortable with the engineering side of data work, dirty formats, inconsistent schemas, pipelines that need to run at three in the morning without anyone watching.

I have thought about what this project would actually require on the ground. The CRITICAL MICROBES database pulls infection test results from laboratories across Sweden. Each lab has its own way of naming pathogens, its own test methods, its own conventions for writing free text descriptions. Before any model can be trained, someone needs to harmonize all of that. That is exactly the kind of problem I solved at SAP, just with clinical data instead of server logs.

The next step is linking this harmonized lab data to population registries, which means joining tables that were never designed to be joined. The result is a heterogeneous graph, patients connected to infections, infections connected to pathogens, pathogens connected to susceptibility profiles, all tracked over time. My master's thesis trained graph attention networks on a molecular interaction graph that combined drug targets, gene expression data, and signed regulatory edges from OmniPath. Different domain, but the same structural challenge: multiple node types, multiple edge types, and the need to learn representations that respect the heterogeneity of the data. I think that experience transfers.

This project sits at the intersection of clinical microbiology, epidemiology, and machine learning. I am not a microbiologist, but I have spent the last two years working with biological knowledge graphs and perturbation data, and I know how to communicate across disciplines. My five years of writing production code and my recent research experience with GNNs put me in a position to contribute from day one.

\textbf{Master's thesis research} During my master's studies, I shifted my focus to computational biology. For my thesis, I built a Counterfactual Attention GNN that predicts drug-induced gene expression changes by combining multi-omics data from LINCS L1000 with biological interaction networks from OmniPath. The project taught me how to connect mathematical models with actual biological findings, and I got hands-on experience evaluating predictive models on real experimental data where the answer is not known in advance.

Thank you for considering my application. I would be very happy to contribute my engineering and research skills to your team.

\renewcommand{\refname}{Selected references}
\bibliographystyle{ieeetr}
\bibliography{references}

Sincerely,

Liuxi Mei

\vspace{0.5cm}
\noindent
\begin{minipage}{\textwidth}
\textbf{References:}

\vspace{0.3cm}
\begin{minipage}[t]{0.48\textwidth}
\textbf{Oleg Sysoev}\\
Senior Associate Professor, Linköping University\\
Department of Computer and Information Science (IDA)\\
Master's Thesis Supervisor\\
Email: \href{mailto:oleg.sysoev@liu.se}{oleg.sysoev@liu.se}
\end{minipage}%
\hfill
\begin{minipage}[t]{0.48\textwidth}
\textbf{Krzysztof Bartoszek}\\
Senior Associate Professor, Linköping University\\
Department of Computer and Information Science (IDA)\\
Teacher and Examiner\\
Email: \href{mailto:krzysztof.bartoszek@liu.se}{krzysztof.bartoszek@liu.se}
\end{minipage}
\end{minipage}

\end{document}
